\documentclass[
  aps,
  pre,
  reprint,
  superscriptaddress,
  nofootinbib
]{revtex4-2}

\usepackage[T1]{fontenc}
\usepackage{lmodern}
\usepackage{amsmath,amssymb,amsfonts,bm,mathtools}
\usepackage{booktabs}
\usepackage{hyperref}

\hypersetup{
  colorlinks=true,
  linkcolor=blue,
  citecolor=blue,
  urlcolor=blue
}

\newcommand{\diag}{\operatorname{diag}}
\newcommand{\Divp}{\operatorname{Div}_{p}}
\newcommand{\Rmass}{R_{\mathrm{mass}}}
\newcommand{\Rsbp}{R_{\mathrm{SBP}}}
\newcommand{\Ddiv}{D_{\mathrm{div}}}

\begin{document}

\title{Split-Parity Mass-Matrix Construction for Folded Spherical SBP Operators}

\author{SphericalSBPOperators Development Team}
\affiliation{Computational Science and Engineering}

\date{\today}

\begin{abstract}
We document an experimental split-mass construction for folded spherical-symmetry
summation-by-parts (SBP) operators.
The method uses distinct norm matrices for parity sectors:
$S$ for even/scalar fields and $V$ for odd/vector fields.
Given a fixed gradient operator $G$, the divergence operator is induced in the interior by
$\Ddiv=-S^{-1}G^{T}V$.
The construction is designed for optimization workflows through explicit structural
constraints on $(S,V)$ by order, monomial exactness constraints on odd regular
test functions, origin constraints, and normalization/asymptotic constraints.
We provide step-by-step implementation details and explicit matrix structures for
orders 2, 4, 6, and 8, and report representative sixth-order diagnostics at grid spacing
$h=1$.
\end{abstract}

\maketitle

\section{Problem Setup}
On a folded half-grid $r_i=ih$, $i=0,\dots,N$, we consider the radial operator
\begin{equation}
  \Divp(u)=\frac{1}{r^p}\partial_r(r^p u),
\end{equation}
with parity split at $r=0$:
even/scalar variable $\Pi$ and odd/vector variable $\Psi$.

The split SBP identity is
\begin{equation}
  S D + G^{T}V = B,
\end{equation}
where $D$ is divergence and $G$ is gradient.
For interior elimination (ignoring $B$ in the elimination step),
\begin{equation}
  \Ddiv \equiv D_{\mathrm{div}} = -S^{-1}G^{T}V.
  \label{eq:ddiv}
\end{equation}

\section{Order-Dependent Structure of \texorpdfstring{$S$}{S} and \texorpdfstring{$V$}{V}}
\subsection{Order 2}
\begin{equation}
  S=\diag(s_0,\dots,s_N), \qquad V=\diag(v_0,\dots,v_N), \qquad V=S.
\end{equation}

\subsection{Order 4 (collocated)}
$S$ is diagonal.
$V$ is symmetric with origin-local off-diagonals only.
In physical node indexing $(0,1,2,\dots)$:
\begin{equation}
  V_{0,1}=V_{1,0}=u_{3/2}, \qquad
  V_{1,2}=V_{2,1}=u_{5/2},
\end{equation}
all other off-diagonals exactly zero.

In the current Julia implementation (1-based matrix indices), this corresponds to
nonzero off-diagonal entries only at $(2,3)$ and $(3,4)$.

\subsection{Order 4 (staggered)}
$S$ is diagonal.
$V$ allows only one origin-local off-diagonal:
\begin{equation}
  V_{0,1}=V_{1,0}=u_1,
\end{equation}
all others zero
(implementation indices $(2,3)$ only).

\subsection{Orders 6 and 8}
$S$ remains diagonal.
$V$ is diagonal in the far field and has a compact symmetric origin block.
Using $N_m=\mathrm{order}/2$ and radius $k=N_m-1$:
\begin{itemize}
  \item order 6: $k=2$, near-origin bandwidth $2k+1=5$;
  \item order 8: $k=3$, near-origin bandwidth $2k+1=7$.
\end{itemize}
Only off-diagonals inside this compact block are allowed.

\section{Optimization-Oriented Constraints}
\subsection{Monomial exactness constraints}
For odd regular monomials
\begin{equation}
  \psi^{(k)}(r)=r^k,\qquad k\in\{1,3,\dots,2N_m-1\},
\end{equation}
the target is
\begin{equation}
  (\partial_r+p/r)\,r^k = (p+k)\,r^{k-1}.
\end{equation}
The linear-in-$(S,V)$ residual form is
\begin{equation}
  \Rmass_i^{(k)} =
  -\bigl(G^T(V\psi^{(k)})\bigr)_i
  - s_i (p+k) r_i^{k-1},
\end{equation}
with hard equality constraints on enforced rows (typically interior rows).

\subsection{Origin constraints}
For centered order 2:
\begin{equation}
  s_1=(1+p)s_0.
\end{equation}
For centered order 4:
\begin{align}
  (1+p)s_0 &= v_1-\tfrac18 u_{3/2}+\tfrac58 u_{5/2}, \\
  v_2 &= v_1+\tfrac{63}{8}u_{3/2}-\tfrac{27}{8}u_{5/2}.
\end{align}

\subsection{Normalization and asymptotic constraints}
Because scaling $S$ and $V$ together leaves Eq.~\eqref{eq:ddiv} unchanged, one
normalization constraint is added (e.g., $s_I=r_I^p$), with optional far-field matching
constraints $s_i\approx r_i^p$ and $v_i\approx r_i^p$ near $r=R$.

\subsection{Positivity and SPD}
Typical parameterization for positivity:
\begin{equation}
  s_i=e^{\sigma_i}, \qquad v_i=e^{\nu_i}.
\end{equation}
For SBP4 centered origin block, SPD can be enforced through a small-block Cholesky
parameterization.

\section{Step-by-Step Procedure}
\begin{enumerate}
  \item Build the folded gradient matrix $G$ on $[0,R]$.
  \item Apply a near-origin gradient repair on selected rows to remove first-column
  coupling while preserving target even-monomial exactness (implemented as an exact
  linear solve on a compact stencil block in the current code path).
  \item Build $(S,V)$ under order-dependent structural constraints.
  \item Evaluate constraint residuals:
  monomial exactness, origin constraints, normalization/far-field constraints.
  \item Form divergence by Eq.~\eqref{eq:ddiv} (or equivalently from the repaired
  SBP relation in the assembled operator path).
  \item Assemble operational $D$ in the solver path (including boundary/origin handling).
  \item Validate:
  $\Rsbp = SD+G^TV-B$, gradient and divergence monomial errors, and quadrature errors.
\end{enumerate}

\section{Current Implementation Notes}
In \texttt{SphericalSBPOperators.jl}, two public helpers were added:
\begin{itemize}
  \item \texttt{construct\_split\_mass\_matrices(...)} for structured $(S,V)$ construction;
  \item \texttt{split\_mass\_constraint\_report(...)} for reporting constraint residuals and SPD/structure checks.
\end{itemize}
Naming convention is fixed as:
$G$ for gradient and $D$ for divergence.

In the default assembly path used for production runs, the folded gradient is
not taken as immutable: a coupled near-origin repair is applied to $G$ before
forming $D$.
This repair solves a local linear system to satisfy even-monomial gradient
conditions together with odd-monomial divergence consistency induced by
$(S,V)$.

\section{Representative Sixth-Order Diagnostics (\texorpdfstring{$h=1$}{h=1})}
Configuration:
order $6$, $N=32$, $R=32$, $p=2$, hence $h=R/N=1$.
Interior diagnostic rows: $8:27$.

\begin{table}[t]
  \caption{Interior and origin errors for derivatives/SBP, and whole-domain quadrature errors.}
  \label{tab:diag6}
  \begin{ruledtabular}
  \begin{tabular}{lll}
    Quantity & Interior error & Origin error \\
    \hline
    $\|\Rsbp\|_\infty$ & $1.42\times 10^{-14}$ & $0$ \\
    grad, $k=0$ & $2.43\times 10^{-17}$ & $0$ \\
    grad, $k=2$ & $2.84\times 10^{-14}$ & $0$ \\
    grad, $k=4$ & $7.28\times 10^{-12}$ & $0$ \\
    grad, $k=6$ & $1.49\times 10^{-8}$ & $0$ \\
    div, $k=1$ & $2.66\times 10^{-15}$ & $0$ \\
    div, $k=3$ & $9.09\times 10^{-13}$ & $0$ \\
  \end{tabular}
  \end{ruledtabular}
\end{table}

Whole-domain quadrature errors:
\begin{align}
  \text{even ($S$):}\;& q=0:\;1.82\times 10^{-12}, \quad q=2:\;9.31\times 10^{-10},\\
  \text{odd ($V$):}\;& q=1:\;8.33\times 10^{-3}, \quad q=3:\;3.89\times 10^{-2}.
\end{align}

\section{Discussion}
The current sixth-order run shows near-machine interior consistency for SBP,
gradient, and divergence diagnostics, and exact origin values in the reported tests.
Odd quadrature errors are comparatively larger, indicating that structural admissibility
alone is not sufficient; optimization of origin-block and far-field mass parameters is
needed to simultaneously satisfy odd quadrature targets and high-order divergence
exactness.

\section{Conclusions}
We established an optimization-ready split-parity mass framework:
\begin{equation}
  (S,V,G)\;\Longrightarrow\;\Ddiv=-S^{-1}G^TV,
\end{equation}
with explicit per-order mass structures, origin constraints, and diagnostics.
The framework is now in place for constrained optimization studies of $(S,V)$
for SBP orders $2/4/6/8$ on collocated and staggered configurations.

\end{document}
